\documentclass[11pt, a4paper]{article}

% ── Packages ──
\usepackage[utf8]{inputenc}
\usepackage[T1]{fontenc}
\usepackage{lmodern}
\usepackage[margin=2.2cm]{geometry}
\usepackage{xcolor}
\usepackage{hyperref}
\usepackage{graphicx}
\usepackage{enumitem}
\usepackage{fancyhdr}
\usepackage{titlesec}
\usepackage{booktabs}
\usepackage{tabularx}
\usepackage{listings}
\usepackage{tcolorbox}
\usepackage{parskip}

% ── Colours ──
\definecolor{jso}{HTML}{6366F1}
\definecolor{darkblue}{HTML}{1E293B}
\definecolor{grey}{HTML}{64748B}
\definecolor{codebg}{HTML}{F8FAFC}
\definecolor{linkblue}{HTML}{4F46E5}

% ── Hyperref ──
\hypersetup{colorlinks=true,linkcolor=jso,urlcolor=linkblue,citecolor=jso}

% ── Header / Footer ──
\pagestyle{fancy}
\fancyhf{}
\fancyhead[L]{\small\color{grey}JSO Candidate Experience Agent}
\fancyhead[R]{\small\color{grey}AariyaTech Corp — Stage 2 Assignment}
\fancyfoot[C]{\small\color{grey}\thepage}
\renewcommand{\headrulewidth}{0.4pt}

% ── Section styling ──
\titleformat{\section}{\Large\bfseries\color{darkblue}}{}{0em}{}[\vspace{-4pt}\color{jso}\rule{\textwidth}{1.5pt}]
\titleformat{\subsection}{\large\bfseries\color{darkblue}}{}{0em}{}
\titleformat{\subsubsection}{\normalsize\bfseries\color{darkblue}}{}{0em}{}

% ── Code listing style ──
\lstset{
  backgroundcolor=\color{codebg},
  basicstyle=\ttfamily\small,
  breaklines=true,
  frame=single,
  rulecolor=\color{gray!30},
  columns=fullflexible,
  keepspaces=true,
  tabsize=2
}

% ────────────────────────────────────────────────────────
\begin{document}

% ═══════ COVER PAGE ═══════
\begin{titlepage}
\centering
\vspace*{3cm}

{\Huge\bfseries\color{jso} JSO}\\[6pt]
{\LARGE\bfseries\color{darkblue} Candidate Experience Agent}\\[20pt]
{\large Stage 2 — Technical Assignment Submission}

\vspace{2cm}

\begin{tabular}{r l}
\textbf{Name} & Mayur Doiphode \\
\textbf{Role Applied For} & Agentic AI Engineer Intern \\
 & (Career Intelligence Systems) \\
\textbf{Company} & AariyaTech Corp Private Limited \\
\textbf{Platform} & JSO (Job Search Optimiser) \\
% \textbf{Date} & \today \\
\end{tabular}

\vspace{2cm}

\begin{tcolorbox}[colback=jso!5, colframe=jso, boxrule=1pt, arc=4pt, width=0.85\textwidth]
\centering
\textbf{GitHub Repository}\\[4pt]
\url{https://github.com/Mayurdoiphode55/Jso.git}

\vspace{10pt}

% \textbf{Live Demo (Vercel)}\\[4pt]
% \url{https://your-app-name.vercel.app}
\end{tcolorbox}

\vfill
{\small\color{grey} All written answers, prototype source code, and deployment instructions are included in the GitHub repository.}
\end{titlepage}

% ═══════ TABLE OF CONTENTS ═══════
\tableofcontents
\newpage

% ═══════════════════════════════════════════════════════
\section*{Links \& Deliverables}
\addcontentsline{toc}{section}{Links \& Deliverables}

\begin{tabularx}{\textwidth}{l X}
\toprule
\textbf{Deliverable} & \textbf{Link / Location} \\
\midrule
GitHub Repository & \url{https://github.com/Mayurdoiphode55/Jso.git} \\
% Live Demo & \url{https://your-app-name.vercel.app} \\
Written Answers & This PDF (Part A \& Part B) \\
Source Code & Full Next.js 14 project in repository \\
Database Schema & \texttt{supabase/migration.sql} in repository \\
\bottomrule
\end{tabularx}

\begin{tcolorbox}[colback=codebg, colframe=gray!30, boxrule=0.5pt, arc=3pt, title={\small\bfseries Demo Credentials}]
\small
\begin{tabular}{lll}
\textbf{Role} & \textbf{Email} & \textbf{Password} \\
Candidate & candidate1@jso.demo & demo123 \\
HR Consultant & consultant1@jso.demo & demo123 \\
Super Admin & admin@jso.demo & demo123 \\
Licensing & licensing@jso.demo & demo123 \\
\end{tabular}
\end{tcolorbox}

\newpage

% ═══════════════════════════════════════════════════════
% PART A
% ═══════════════════════════════════════════════════════
\section*{\LARGE PART A — WRITTEN ANSWERS}
\addcontentsline{toc}{section}{PART A — Written Answers}

% ── Section 1 ──
\section{Section 1: Why Agentic JSO}

\subsection{Q1: Why does the JSO platform require AI Agents in Phase 2?}

The JSO platform's Phase 1 architecture is built entirely on manual processes — human HR consultants conducting one-on-one sessions, manual scheduling, and subjective quality assessments. While this approach works at a small scale, it fundamentally cannot scale to serve the global job-seeking population that JSO's mission targets. When the platform grows from hundreds to tens of thousands of concurrent users across multiple time zones and languages, human-only operations become a bottleneck rather than a backbone.

AI agents solve this by operating 24/7 without fatigue, cognitive bias, or capacity constraints. A single AI agent can simultaneously process feedback from thousands of consultation sessions, detect patterns that human administrators would miss, and generate actionable insights in real time. For instance, a human admin reviewing consultation quality might read 10--15 feedback forms per day; an AI agent can analyse thousands within seconds, catching subtle sentiment shifts that indicate declining consultation quality before they become systemic problems.

JSO's core mission — solving the global job search problem — is only achievable at scale with AI. Personalised career guidance for thousands of users simultaneously requires intelligent automation that can adapt to each candidate's context, language, and needs. AI agents don't replace human consultants; they amplify them, handling repetitive analytical tasks so consultants can focus on what they do best — empathetic, personalised career guidance.

From a \textbf{Governance} perspective, AI agents bring unprecedented transparency and accountability to platform quality. Every agent decision is logged, every analysis is traceable, and every alert has a clear audit trail. This creates a governance framework where platform quality is not just aspirational but measurable and enforceable — something that pure manual operations can never achieve consistently at scale.

\subsection{Q2: What inefficiencies exist in the current Phase 1 system?}

The current Phase 1 system has several critical inefficiencies that directly impact candidate outcomes and platform quality:

\textbf{No automated feedback collection.} After a consultation session ends, there is no systematic mechanism to capture the candidate's experience. Whether a session was transformative or deeply unsatisfying, the platform has zero visibility.

\textbf{No real-time quality monitoring.} Consultant performance is essentially a black box. There are no metrics, no dashboards, and no alerts. A consultant could deliver consistently poor sessions for weeks before anyone notices.

\textbf{Zero admin visibility into consultation quality.} Super Admins and platform managers operate without data. They cannot identify which consultants are excelling, which need support, or which sessions are generating complaints.

\textbf{Manual, subjective consultant-candidate matching.} Without data-driven insights into consultant strengths and candidate needs, matching is essentially random or based on availability alone.

\textbf{No automated follow-up system.} After a session, candidates receive no check-in, no resources, and no nudge to book follow-ups.

From a \textbf{Workers} pillar perspective, this lack of structured feedback is equally harmful to consultants. Without systematic performance data, HR consultants have no mechanism for professional growth. They cannot see what they're doing well, what candidates value, or where they need to improve. Feedback is a gift, and the current system denies consultants this essential growth tool.

\subsection{Q3: How can AI agents improve user experience and platform efficiency?}

AI agents can transform the JSO platform from a reactive, manually-operated system into a proactive, intelligence-driven career platform:

\textbf{Automated post-session satisfaction surveys triggered instantly.} Within minutes of a consultation ending, the Candidate Experience Agent sends a short, focused survey. This captures feedback while the experience is fresh, improving both response rates and feedback quality.

\textbf{Sentiment analysis on feedback text to detect frustration or praise.} Using Groq's LLaMA 3.1 model, the agent analyses free-text comments to detect nuanced emotions and extracts specific issues like ``salary question ignored'' or ``excellent resume advice.''

\textbf{Real-time alerts to admins when quality drops.} The moment a session receives a satisfaction score below the threshold, the Super Admin dashboard receives an automatic alert.

\textbf{Personalised nudges to candidates for follow-up bookings.} Based on session outcomes and satisfaction data, the agent can recommend follow-up sessions or tailored resources.

\textbf{Trend reports helping consultants improve.} Weekly and monthly performance trends give consultants a clear picture of their trajectory.

From a \textbf{Community} pillar perspective, AI agents enable multilingual survey support, ensuring candidates from diverse linguistic backgrounds can provide feedback in their preferred language.

% ── Section 2 ──
\section{Section 2: Agent Design}

\subsection{Q1: What type of AI agent should be built?}

The Candidate Experience Agent (CEA) should be designed as a \textbf{Reactive + Goal-Oriented AI Agent}, combining event-driven responsiveness with strategic optimisation toward platform quality goals.

\textbf{Reactive Component:} The agent responds to specific events in real time. When a consultation session status changes to `completed' in the database, the agent immediately triggers the survey workflow. When a candidate submits feedback, the agent instantly processes it through the AI analysis pipeline. When a satisfaction score falls below the threshold, the agent creates an alert without delay.

\textbf{Goal-Oriented Component:} Beyond simply reacting to events, the CEA works toward a defined goal: maximising overall candidate satisfaction scores across the platform. It generates improvement suggestions for consultants, identifies systemic patterns, and provides the data foundation for strategic decisions.

The CEA uses \textbf{Groq's LLaMA 3.1 8B Instant model} for natural language understanding, enabling it to parse free-text feedback with nuance that simple keyword matching cannot achieve.

From a \textbf{Governance} perspective, every decision the agent makes is fully logged and auditable. The AI analysis results are stored alongside the raw feedback, creating a transparent chain from candidate comment to platform action.

\subsection{Q2: What tasks will this agent automate?}

The Candidate Experience Agent automates six core tasks:

\begin{enumerate}[leftmargin=2em]
\item \textbf{Auto-send satisfaction survey after every session.} When a session is marked `completed,' the agent triggers a short survey — a star rating (1--5) plus an optional free-text comment.
\item \textbf{Analyse free-text feedback with AI.} The candidate's comment is anonymised (names and emails stripped) and sent to Groq's LLaMA 3.1 8B model for structured analysis.
\item \textbf{Calculate satisfaction score (1--10).} The agent combines the star rating with the AI's sentiment analysis using a blended formula.
\item \textbf{Flag low-score sessions and create alerts.} Any session receiving a score below 6/10 is automatically flagged with an alert containing key issues.
\item \textbf{Generate weekly trend reports.} The agent aggregates satisfaction data to produce trend visualisations.
\item \textbf{Recommend improvement areas to consultants.} Based on patterns in flagged sessions, the agent generates targeted improvement suggestions.
\end{enumerate}

\subsection{Q3: How will the agent interact with the existing dashboard?}

The CEA integrates with all four JSO dashboards:

\textbf{User (Candidate) Dashboard:} Post-session survey popup with interactive star rating and privacy notice. Personal satisfaction history with past ratings.

\textbf{HR Consultant Dashboard:} Real-time satisfaction score widget (colour-coded), trend line chart, recent feedback with AI sentiment badges and key issues. From the \textbf{Workers} pillar, feedback is framed as ``Here's how you can grow'' rather than punitive.

\textbf{Super Admin Dashboard:} Platform-wide metrics, satisfaction trend chart, consultant leaderboard, active alerts panel with ``Mark Resolved'' actions. Full audit trail from the \textbf{Governance} pillar.

\textbf{Licensing Dashboard:} Aggregate-only metrics, compliance summaries, quarterly trend visualisation. No individual candidate data is visible.

% ── Section 3 ──
\section{Section 3: Problem Solving}

\subsection{Q1: What specific problem does this agent solve?}

The fundamental problem is the complete absence of a systematic mechanism to evaluate, monitor, and improve the quality of consultation sessions. In Phase 1, consultation quality is invisible — there are no metrics, no feedback loops, and no automated quality assurance.

\textbf{Poor consultants go undetected and uncorrected.} Without detection, there can be no correction, and without correction, quality deteriorates unchecked.

\textbf{Candidates leave with unresolved issues and don't return.} Silent churn is the most dangerous kind because it's entirely invisible to platform operators.

\textbf{JSO cannot improve service quality systematically.} Without quantitative data, platform improvements are based on guesswork.

From the \textbf{Customers} pillar perspective, candidates currently have no voice in platform quality. The CEA solves this by giving candidates a structured channel to share their experiences and ensuring that every piece of feedback drives measurable action.

\subsection{Q2: Provide a real scenario showing the agent's impact}

\textbf{Scenario: Priya's Consultation Experience}

Priya Sharma books a consultation session. The consultant rushes through her questions, provides generic advice, and dismisses her salary negotiation question. Priya leaves frustrated.

\textbf{In Phase 1:} Nothing happens. The consultant continues the same way. Priya considers cancelling.

\textbf{In Phase 2 (with CEA):} Five minutes after the session, the CEA sends a satisfaction survey. Priya rates it 2/5 and writes her feedback. The CEA strips PII and sends it to Groq:

\begin{lstlisting}
{
  "sentiment": "negative",
  "key_issues": ["generic advice", "salary question dismissed", "felt rushed"],
  "sentiment_score": -0.78,
  "summary": "Candidate expressed strong dissatisfaction with rushed, generic consultation."
}
\end{lstlisting}

The CEA calculates a satisfaction score of \textbf{3/10}. Since 3 < 6, the session is flagged and an alert is created. The consultant's dashboard reflects the new score with improvement suggestions. The Super Admin receives the alert.

From a \textbf{Sustainability} pillar perspective, this systematic quality loop helps underserved job seekers get progressively better guidance.

% ── Section 4 ──
\section{Section 4: Dashboard Integration}

\subsection{A. User (Candidate) Dashboard}

\begin{itemize}[leftmargin=2em]
\item \textbf{Post-consultation survey popup} — 5-star rating + optional comment + privacy notice
\item \textbf{Personal satisfaction history} — chronological list of past feedback with AI summaries
\item \textbf{Session management} — upcoming/past sessions with ``Give Feedback'' buttons
\end{itemize}

\subsection{B. HR Consultant Dashboard}

\begin{itemize}[leftmargin=2em]
\item \textbf{Real-time satisfaction score widget} — colour-coded (green $\geq$ 7, yellow 5--6, red < 5)
\item \textbf{Flagged sessions with AI improvement suggestions}
\item \textbf{Weekly performance trend graph} — line chart of score trajectory
\end{itemize}

From the \textbf{Workers} pillar, the dashboard language is carefully chosen: ``Here's how you can grow'' rather than ``Here's what you did wrong.''

\subsection{C. Super Admin Dashboard}

\begin{itemize}[leftmargin=2em]
\item \textbf{Platform-wide satisfaction trend graph} — area chart with monthly/weekly view
\item \textbf{Unresolved alerts list} — with details and ``Mark Resolved'' actions
\item \textbf{Consultant leaderboard} — ranked by satisfaction score
\end{itemize}

\subsection{D. Licensing Dashboard}

\begin{itemize}[leftmargin=2em]
\item \textbf{Aggregate quality metrics} — no individual candidate data visible
\item \textbf{Compliance-ready reports} — export for regulatory submissions
\item \textbf{Quarterly trend visualisation} — bar chart of monthly scores
\end{itemize}

% ── Section 5 ──
\section{Section 5: Technical Architecture}

\subsection{Components}

\begin{tabularx}{\textwidth}{l l X}
\toprule
\textbf{Component} & \textbf{Technology} & \textbf{Purpose} \\
\midrule
Frontend & Next.js 14 + React + Tailwind & Server-rendered UI with App Router \\
Backend & Next.js API Routes & Stateless serverless endpoints \\
AI Engine & Groq API — LLaMA 3.1 8B & Sentiment analysis + key issue extraction \\
Database & Supabase PostgreSQL & Structured data, Auth, Realtime \\
Hosting & Vercel & Serverless deployment \\
\bottomrule
\end{tabularx}

\subsection{System Architecture Diagram}

\begin{center}
\includegraphics[width=\textwidth]{architecture.png}
\end{center}
{\small\itshape Figure: End-to-end system architecture — candidate submits feedback through the frontend, API routes process and analyse via Groq AI, results are stored in Supabase, and Realtime pushes updates to all dashboards.}

\subsection{End-to-End Flow}

\begin{enumerate}[leftmargin=2em]
\item Session marked `completed' in Supabase
\item Survey popup displayed to candidate
\item Candidate submits star rating + comment
\item PII stripped from comment text
\item Anonymised text sent to Groq LLaMA 3.1
\item AI returns sentiment, key issues, score, summary
\item Satisfaction score calculated (blended formula)
\item Survey record updated in Supabase
\item Score < 6 $\rightarrow$ Alert created automatically
\item Supabase Realtime pushes updates to all dashboards
\end{enumerate}

From the \textbf{Environment} pillar, this serverless architecture means compute resources only activate when needed. No always-on servers consuming energy.

\subsection{API Endpoints}

\begin{tabularx}{\textwidth}{l l X}
\toprule
\textbf{Endpoint} & \textbf{Method} & \textbf{Purpose} \\
\midrule
\texttt{/api/surveys/create} & POST & Creates a new survey after session ends \\
\texttt{/api/surveys/analyse} & POST & Calls Groq API and stores AI analysis \\
\texttt{/api/surveys/consultant/[id]} & GET & Returns consultant's satisfaction data \\
\texttt{/api/alerts} & GET & Returns unresolved alerts for Super Admin \\
\texttt{/api/alerts} & PATCH & Marks an alert as resolved \\
\bottomrule
\end{tabularx}

% ── Section 6 ──
\section{Section 6: Integration With Phase 1}

\subsection{Data Flow}

\begin{lstlisting}
Candidate submits survey in React UI
  -> Next.js API route (/api/surveys/create)
  -> PII stripped from comment text
  -> Groq API (LLaMA 3.1 8B analysis)
  -> Results stored in Supabase (surveys table)
  -> Score < 6? -> Alert created (alerts table)
  -> Supabase Realtime broadcasts changes
  -> All 4 dashboards update simultaneously
\end{lstlisting}

\subsection{Security (MCP/APA Compliance)}

\textbf{Authentication:} All API routes are protected with JWT tokens issued by Supabase Auth.

\textbf{Row Level Security (RLS):} Every Supabase table has RLS policies enforcing role-based access.

\textbf{Data Anonymisation:} PII-stripping function removes names, emails, and phone numbers before any data reaches the Groq API.

\textbf{Transport Security:} All communications use HTTPS. API keys are stored as server-side environment variables.

From the \textbf{Customers} pillar, candidate data is never exposed raw to AI systems. The anonymisation layer is structurally enforced in the API pipeline.

% ── Section 7 ──
\section{Section 7: Build Timeline}

\begin{tabularx}{\textwidth}{l X}
\toprule
\textbf{Week} & \textbf{Activities} \\
\midrule
Week 1 & Architecture design, API contracts, Supabase schema, RLS policies \\
Week 2--3 & Core agent: survey system, Groq integration, scoring engine, alerts, PII stripping \\
Week 4 & Dashboard widgets for all 4 roles, Recharts integration, Realtime subscriptions \\
Week 5 & Testing: unit, integration, security audit, cross-browser \\
Week 6 & Vercel deployment, monitoring, documentation, demo data, final QA \\
\bottomrule
\end{tabularx}

\newpage

% ═══════════════════════════════════════════════════════
% PART B
% ═══════════════════════════════════════════════════════
\section*{\LARGE PART B — DESIGN DOCUMENT}
\addcontentsline{toc}{section}{PART B — CEA Design Document}

\section{Problem Statement}

The JSO platform currently operates without any mechanism to evaluate consultation quality. When a session ends, no data is captured, no analysis is performed, and no action is taken. Poorly performing consultants go undetected, candidates silently disengage, and platform administrators make decisions based on intuition rather than evidence. The Candidate Experience Agent addresses this by creating an automated, AI-powered feedback loop.

\section{Agent Overview}

\begin{tabularx}{\textwidth}{l X}
\toprule
\textbf{Attribute} & \textbf{Detail} \\
\midrule
Name & Candidate Experience Agent (CEA) \\
Type & Reactive + Goal-Oriented AI Agent \\
Purpose & Automate satisfaction measurement, analyse feedback, flag issues, drive improvement \\
AI Model & Groq LLaMA 3.1 8B Instant \\
Primary Metric & Platform-wide average satisfaction score (target $\geq$ 7/10) \\
Data Store & Supabase PostgreSQL (surveys, alerts tables) \\
\bottomrule
\end{tabularx}

\section{Step-by-Step Flow}

\begin{enumerate}[leftmargin=2em]
\item Session marked `completed' in Supabase
\item Candidate dashboard detects completed session
\item Survey popup displayed (star rating + comment)
\item Candidate submits feedback
\item API route receives submission
\item PII stripped from comment text
\item Anonymised text sent to Groq LLaMA 3.1 8B
\item AI returns: sentiment, key\_issues, score, summary
\item Satisfaction score calculated (blended formula)
\item Survey record updated in Supabase
\item Score < 6 $\rightarrow$ Alert created automatically
\item Supabase Realtime pushes updates to all dashboards
\end{enumerate}

\section{Scoring System}

\subsection{Formula}

\begin{lstlisting}
satisfaction_score = Math.round(
  ((star_rating / 5) + ((sentiment_score + 1) / 2)) / 2 * 10
)
\end{lstlisting}

\subsection{Score Interpretation}

\begin{tabularx}{\textwidth}{l l l X}
\toprule
\textbf{Score} & \textbf{Classification} & \textbf{Colour} & \textbf{Action} \\
\midrule
8--10 & Excellent & Green & No action needed \\
7 & Good & Green & No action needed \\
5--6 & Needs Attention & Yellow & Monitor trend \\
1--4 & Poor & Red & Alert generated, intervention required \\
\bottomrule
\end{tabularx}

\section{Alert System}

An alert is created automatically when a survey's satisfaction score is below 6/10. Alerts contain the consultant ID, score, key issues, and are visible to Super Admins who can mark them as resolved.

\subsection{Alert Lifecycle}

\begin{enumerate}[leftmargin=2em]
\item \textbf{Created} — Automatically when score < 6
\item \textbf{Visible} — Appears on Super Admin dashboard via Realtime
\item \textbf{Reviewed} — Admin reads details and key issues
\item \textbf{Resolved} — Admin clicks ``Mark Resolved'' after taking action
\item \textbf{Archived} — Resolved alerts move to historical record
\end{enumerate}

\section{AI Analysis Logic}

\subsection{Prompt Structure}

\begin{lstlisting}
System: You are a sentiment analysis engine for a career consultation 
platform. Analyse the candidate feedback and return ONLY a JSON object:
{
  "sentiment": "positive" | "neutral" | "negative",
  "key_issues": ["issue1", "issue2"],
  "sentiment_score": <number between -1.0 and 1.0>,
  "summary": "<one sentence summary>"
}

User: [anonymised candidate feedback text]
\end{lstlisting}

\subsection{PII Stripping}

Before sending to Groq, all personally identifiable information is removed: names $\rightarrow$ [CANDIDATE], emails $\rightarrow$ [EMAIL], phone numbers $\rightarrow$ [PHONE].

\section{Privacy \& Ethics}

\begin{enumerate}[leftmargin=2em]
\item \textbf{PII Stripping} — All identifiable information removed before AI analysis
\item \textbf{Row Level Security} — Supabase RLS enforces role-based data access
\item \textbf{JWT Authentication} — Every API request is authenticated
\item \textbf{Anonymised Consultant Views} — Candidate names shown as ``Candidate A, B, C''
\item \textbf{No Raw Data in Licensing} — Only aggregate metrics, no individual responses
\item \textbf{HTTPS Only} — All data encrypted in transit
\item \textbf{Consent-Based Feedback} — Clear privacy notice on every survey form
\end{enumerate}

\section{Part C Pillars Integration}

\begin{tabularx}{\textwidth}{l X}
\toprule
\textbf{Pillar} & \textbf{How CEA Addresses It} \\
\midrule
Governance & All agent decisions logged and auditable. Full transparency in scoring, alerting, AI analysis. RLS policies enforce data access controls. \\
Workers & Consultant feedback framed as professional growth, not surveillance. Improvement suggestions are constructive and specific. \\
Community & Survey design supports multilingual candidates. Anonymisation ensures feedback is judged on content, not identity. \\
Environment & Serverless architecture (Vercel + Supabase) activates compute only on demand. Groq LPU minimises carbon footprint per AI call. \\
Customers & PII stripping is structurally enforced. Candidates receive clear privacy notices. Feedback directly drives service improvements. \\
Sustainability & Systematic quality improvement creates a virtuous cycle: better feedback $\rightarrow$ better consultants $\rightarrow$ better outcomes $\rightarrow$ higher platform value. \\
\bottomrule
\end{tabularx}

\vspace{1cm}

\begin{center}
\color{grey}\rule{0.5\textwidth}{0.5pt}\\[10pt]
{\small\itshape End of Assignment Submission}
\end{center}

\end{document}
